\begin{afterword}
\summaryhead

The essay warns that teaching people about cognitive biases can make them reason worse. In its opening story the author's mother hears about expert overconfidence from the author and is delighted to have a new way to dismiss experts.

The evidence is a study by Taber and Lodge, which found that people judge arguments on their side more kindly than arguments on the other, and that politically knowledgeable people do this more, because they have more material for counterargument. The author concludes that knowledge can hurt someone who is already irrational, and that knowing about biases can become ammunition against any conclusion one dislikes.

Two anecdotes follow, about one man. He used overconfidence research to dismiss experts; then, taught about sophisticated arguers, he accused the author of being one. The author now opens any talk on biases with a long warning about motivated skepticism before mentioning calibration.

\responsehead

The title makes a claim that could be tested: that learning about biases makes people reason worse. The essay does not test it. Its one study concerns political knowledge, facts about politics, and says nothing about people who have studied biases. Everything else is anecdote, told by the author about unnamed people, in which the author is the one who understood and the others misused what they learned.

The essay's opening also misstates the finding it uses. The author's mother is told that experts at 99 per cent confidence are right about 70 per cent of the time. The best-known result with those numbers came from students answering general-knowledge questions, not from experts. Some experts, such as weather forecasters, are well calibrated.

Throughout, the post does what it warns against. It warns against using bias-talk on others, then asks the reader, ``You can think of people who fit this description, right?'' It warns that the notion of a ``sophisticated arguer'' is deadly, then reports a talk that spent thirty minutes on sophisticated arguers. At each step the remedy for bias-talk is more bias-talk, delivered by the author, aimed at other people.

The essay's proposed fix is also untested. Putting disconfirmation bias before calibration in a lecture changes the order of the warnings. It does nothing about the problem the anecdotes show, which is that listeners who hear the warnings use them on others. The author gives no evidence that any audience reasoned better after such talks.

In short: a warning that people use knowledge of biases to dismiss others, which opens by misstating a bias finding, invites the reader to apply its label to others, and proposes more of the same teaching as the remedy.
\end{afterword}
